% L'ambiente "abstract" e' definito dalla classe: apre una pagina nuova,
% a piena larghezza, con intestazione e filetto.
%
% Stile: abstract descrittivo sul modello della tesi Scattolini (relatore
% Callisto, a.a. 2022/23), indicata dal professore come esempio di testo:
% forma impersonale, capoversi con ruolo fisso (obiettivo, studio, metodo,
% sviluppo, conclusione), nessun numero. Testo dell'autore (16/09/2026) con tre
% correzioni di sostanza: il PIR e' quello installato nel nodo dispiegato, non
% un'idea di partenza (paper DIPME, Tab. 1); il radar rileva "in tutti gli
% scenari di confronto", non "in tutte le condizioni" (a 120 gradi e a 7 m non
% rileva, cap. 4); nella campagna si misura il respiro, l'indice di vitalita' e'
% la proposta del cap. 7. I valori quantitativi sono nel cap. 4 e nel cap. 8.
\begin{abstract}

La tesi analizza e confronta diverse tecnologie per il rilevamento della presenza
umana in ambienti indoor, nell'ambito del progetto DIPME. Il
progetto prevede l'integrazione di sensori IoT (\emph{Internet of Things}) negli
arredi scolastici, con l'obiettivo di segnalare, in seguito a un evento sismico,
l'eventuale presenza di persone che potrebbero essersi rifugiate sotto di essi.

Il nodo attualmente installato negli arredi affida il rilevamento della presenza a
un sensore PIR (\emph{Passive InfraRed}). Tuttavia, a causa del suo principio di
funzionamento, il PIR è in grado di rilevare il movimento, ma non una persona
immobile. Il lavoro valuta quindi l'impiego di un radar \emph{mmWave} FMCW
(\emph{Frequency Modulated Continuous Wave}) a 24~GHz, per verificare se possa
superare questo limite e quali siano i compromessi legati al suo utilizzo.

Nella prima parte vengono studiati il sensore PIR HC-SR501 e i due moduli radar
HLK-LD2410B e HLK-LD2420, con particolare attenzione ai loro principi di
funzionamento, ai protocolli di comunicazione, ai parametri disponibili e ai limiti
riportati nella documentazione dei costruttori. Per il confronto è stata realizzata
una piattaforma di acquisizione basata su ESP32, in grado di campionare i sensori
in modo sincrono e di accedere, nel caso dei radar, all'energia riflessa per
ciascuna cella di distanza. A supporto della campagna è stata progettata e
implementata un'interfaccia web, servita direttamente dall'ESP32, per la
visualizzazione dei dati in tempo reale e la loro esportazione nello stesso formato.

La piattaforma è stata utilizzata per una campagna sperimentale che ha preso in
esame la distanza di rilevamento, le latenze, i falsi positivi, l'attenuazione del
segnale in presenza di ostacoli, la copertura angolare e il rilevamento del respiro.
I risultati confermano che il PIR risponde al movimento e non alla persona ferma,
mentre il radar l'ha rilevata anche da immobile, in tutti gli scenari di confronto.
Anche il secondo radar ha confermato questo comportamento, pur presentando alcuni
limiti propri.

Sulle grandezze graduate fornite dal radar è stato inoltre sviluppato un
\emph{indice di vitalità} calcolato direttamente a bordo, che, oltre a indicare la
presenza di una persona, fornisce una stima del suo livello di movimento. Il
consumo energetico dei sensori è stato infine valutato sulla base dei dati forniti
dai rispettivi costruttori.

Il lavoro si conclude con la proposta di un'architettura ibrida in cui il PIR,
grazie al suo basso consumo energetico, viene utilizzato per determinare quando
attivare il radar.

\end{abstract}
